\documentclass{book}
\usepackage{graphicx}                              %for PNG images
\usepackage[linkbordercolor={1.0 1.0 0.0}]{hyperref} %for \url tag
\usepackage{color}                                 %for defining custom colors
\usepackage{framed}                                %for shaded and framed paragraphs
\usepackage{textcomp}                              %for various symbols, e.g. Registered Mark
\usepackage{makeidx}                               %for index page
\usepackage{geometry}                              %for defining page size
\usepackage{longtable}                             %for breaking tables
%
\geometry{verbose,a4paper,tmargin=2.5cm,bmargin=2.5cm,lmargin=2.5cm,rmargin=2cm}
\makeindex
\hypersetup{
  pdfauthor = {Mattias Ellert},
  pdftitle = {The ARC Data Manager Component (DMC) Implementation Guide},
  pdfsubject = {DMC Implemenatation Guide},
  pdfkeywords = {Grid,NorduGrid,ARC,DMC},
  pdfcreator = {PDFLaTeX with hyperref package},
  pdfproducer = {PDFLaTeX}
}
%
\bibliographystyle{IEEEtran}
%
\def\efill{\hfill\nopagebreak}%
\hyphenation{Nordu-Grid}
\setlength{\parindent}{0cm}
\setlength{\FrameRule}{1pt}
\setlength{\FrameSep}{8pt}
\addtolength{\parskip}{5pt}
\renewcommand{\thefootnote}{\fnsymbol{footnote}}
\renewcommand{\arraystretch}{1.3}
\newcommand{\dothis}{\colorbox{shadecolor}}
\definecolor{shadecolor}{rgb}{1,1,0.6}
\definecolor{salmon}{rgb}{1,0.9,1}
\definecolor{bordeaux}{rgb}{0.75,0.,0.}
\definecolor{cyan}{rgb}{0,1,1}
%
\begin{document}

\def\today{\number\day/\number\month/\number\year}

\begin{titlepage}

  \begin{tabular}{rl}
    \resizebox*{3cm}{!}{\includegraphics{ng-logo.png}}
    &\parbox[b]{2cm}{\textbf \it {\hspace*{-1.5cm}NORDUGRID\vspace*{0.5cm}}}
  \end{tabular}

  \hrulefill

  {\raggedleft NORDUGRID-TECH-23\par}

  {\raggedleft \today\par}

  \vspace*{2cm}

  {\centering \textsc{\Large ARC Data Manager Component (DMC)}\Large \par}

  \vspace*{0.5cm}

  {\centering \textit{\large Implementation Guide}\large \par}

  \vspace*{1.5cm}

  {\centering \large Mattias Ellert\footnote{mattias.ellert@fysast.uu.se} \large \par}

\end{titlepage}

\tableofcontents
%\listoffigures
\newpage
\chapter{Introduction}
\label{sec:intro}

The libarcdata2 library in the new Advanced Resource Connector (ARC)
is based on libarcdata library in the old ARC. Unlike its predecessor
the libarcdata2 library takes advantage of the plugin mechanism of the
ARC Loader. The core libarcdata2 library does not introduce any
additional external dependencies. Plugins for specific data access
protocols can however have various external dependencies. This
separation of external dependencies from the core library helps reduce
the minimum set of requirements for ARC while allowing the support for
additional access protocols requiring special dependencies to be
installed by those who need it.

The plugins used by the libarcdata2 library are called Data Manager
Components (DMCs). Each DMC should provide a class that inherits from
the DMC base class. Like other ARC Loader plugins the DMCs should
provide a descriptor that maps the name of the plugin to an function
that creates an instance of the DMC. In order for the default DMC
client configuration to automatically find the DMC the name of the DMC
in the descriptor and the name of the loadable module containing the
DMC should match. A DMC called X should reside in a module called
libdmcX.

The constructor and destructor of a DMC should call the Register and
Unregister method of respectively.

The most central class in the arcdata2 library is the DataPoint
class. This class describes a source or destination of a file
transfer.

A DMC should provide an iGetDataPoint method that returns a pointer to
an instance of a specialisation of the DataPoint class when called
with a URL. If the DMC does not support the provided URL a NULL
pointer should be returned instead.

The libarcdata2 library contains a wrapper class for the DataPoint
class called DataHandle. Calling the constructor of the DataHandle
class with a URL will trigger the iGetDataPoint class of all
registered DMCs to be called until a non-NULL pointer is returned.
By dereferencing the DataHandle object the underlying DataPoint can be
accessed.

Two partial specialisations of the DataPoint class are available in
the libarcdata2 library. These are the DataPointIndex and
DataPointDirect classes. The DataPointIndex class implements common
functionality for DMCs that support file catalogues, while the
DataPointDirect implement common functionality for DMCs that support
direct file access. The specialisation of the DataPoint class used by
a DMC should normally inherit from one of these classes rather than
from the DataPoint base class directly.

\chapter{The DMC class}
\label{sec:DMC}

The DMC class describes a loadable module for arcdata2 library. A
typical specialization of this class for a custom DMC plugin would
look like this:

\begin{verbatim}
#ifndef __ARC_DMCCUSTOM_H__
#define __ARC_DMCCUSTOM_H__

#include <arc/data/DMC.h>

namespace Arc {

  class DMCCustom : public DMC {
   public:
    DMCCustom(Config *cfg);
    ~DMCCustom();
    static DMC* Instance(Config *cfg, ChainContext *ctx);
    DataPoint* iGetDataPoint(const URL& url);
   protected:
    static Logger logger;
  };

} // namespace Arc

#endif // __ARC_DMCCUSTOM_H__
\end{verbatim}

The corresponding implementation should look like this:

\begin{verbatim}
#include <arc/Logger.h>
#include <arc/URL.h>
#include <arc/loader/DMCLoader.h>

#include "DMCCustom.h"
#include "DataPointCustom.h"

namespace Arc {

  Logger DMCCustom::logger(DMC::logger, "Custom");

  DMCCustom::DMCCustom(Config *cfg) : DMC(cfg) {
    Register(this);
  }

  DMCCustom::~DMCCustom() {
    Unregister(this);
  }

  DMC* DMCCustom::Instance(Arc::Config *cfg, Arc::ChainContext*) {
    return new DMCCustom(cfg);
  }

  DataPoint* DMCCustom::iGetDataPoint(const URL& url) {
    if (url.Protocol() != "custom") return NULL;
    return new DataPointCustom(url);
  }

} // namespace Arc

dmc_descriptors ARC_DMC_LOADER = {
  { "custom", 0, &Arc::DMCCustom::Instance },
  { NULL, 0, NULL }
};
\end{verbatim}

The DataPointCustom class in the above example is the specialization
of the DataPoint class for this DMC (see chapter \ref{sec:DataPoint}).

The DMC can support more than one protocol in which case the
iGetDataPoint should be changed accordingly.

\chapter{The DataPoint class}
\label{sec:DataPoint}

The DataPoint class is the most central class in the arcdata2
library. This is an abstract base class that must be specialized for
each DMC. Two partial specialisations of the DataPoint class are
available in the libarcdata2 library. These are the DataPointIndex and
DataPointDirect classes. The DataPointIndex class implements common
functionality for DMCs that support file catalogues, while the
DataPointDirect implement common functionality for DMCs that support
direct file access. The specialisation of the DataPoint class used by
a DMC should normally inherit from one of these classes rather than
from the DataPoint base class directly.

\section{DataPointDirect}

To create a DataPoint class for a direct access protocol, the
following methods should be provided:

\begin{verbatim}
namespace Arc {

  class DataPointCustomDirect : public DataPointDirect {
   public:
    DataPointCustomDirect(const URL& url);
    virtual ~DataPointCustomDirect();
    virtual DataStatus StartReading(DataBuffer& buffer);
    virtual DataStatus StartWriting(DataBuffer& buffer,
                                    DataCallback *space_cb = NULL);
    virtual DataStatus StopReading();
    virtual DataStatus StopWriting();
    virtual DataStatus Check();
    virtual DataStatus Remove();
    virtual DataStatus ListFiles(std::list<FileInfo>& files,
                                 bool long_list = false, bool resolve = false);
  };

} // namespace Arc
\end{verbatim}

\section{DataPointIndex}

To create a DataPoint class for an index protocol, the following
methods should be provided:

\begin{verbatim}
namespace Arc {

  class DataPointCustom2 : public DataPointIndex {
   public:
    DataPointCustom2(const URL& url);
    ~DataPointCustom2() {};
    virtual DataStatus Resolve(bool source);
    virtual DataStatus PreRegister(bool replication, bool force = false);
    virtual DataStatus PostRegister(bool replication);
    virtual DataStatus PreUnregister(bool replication);
    virtual DataStatus Unregister(bool all);
    virtual DataStatus ListFiles(std::list<FileInfo>& files,
                                 bool long_list = false, bool resolve = false);
  };

} // namespace Arc
\end{verbatim}

%\bibliography{grid}
%\printindex

\end{document}
